\chapter{Manifold-merged entanglement entropy}
\label{ch:manifold-entropy}

\section{Purpose and scope}

The active-space criterion of this package ranks orbitals by the
four-state entanglement entropy of their one-orbital reduced density
matrix.  The entropy is defined for one state, and the release
computed it for one state.  A spin multiplet is not one state: it is a
manifold of $2S+1$ degenerate components related by spin projection,
and the component that a single-reference calculation returns is a
poor representative of the manifold for the purpose of this
criterion.  This chapter defines the merge that fixes the
representative, gives the derivation, documents the API, and describes
the verification gate at the robust level that the merge now feeds.

\section{Why one state is not enough}
\label{sec:manifold-why}

\subsection{The four occupation weights}

For one orbital, the one-orbital reduced density matrix in the
occupation basis $\{|0\rangle,|\uparrow\rangle,|\downarrow\rangle,
|\uparrow\downarrow\rangle\}$ is diagonal, and its four weights are
\begin{equation}
  w = \bigl(\,1-n_{\alpha}-n_{\beta}+p_{2},\;\;
             n_{\alpha}-p_{2},\;\;
             n_{\beta}-p_{2},\;\;
             p_{2}\,\bigr),
  \label{eq:four-state-weights}
\end{equation}
where $n_{\alpha}$, $n_{\beta}$ are the spin-resolved occupations of
the orbital and $p_{2}$ is the probability of finding it doubly
occupied.  The entropy is
\begin{equation}
  s_{1} = -\sum_{k=1}^{4} w_{k}\ln w_{k},
  \qquad 0 \le s_{1} \le \ln 4 .
  \label{eq:four-state-entropy}
\end{equation}
An orbital that is always empty, always singly occupied with a fixed
spin, or always doubly occupied has $s_{1}=0$; an orbital whose
occupation fluctuates strongly approaches $\ln 4$.

\subsection{The $M_S=+S$ component of a high-spin determinant}

Consider the component of largest spin projection of a pure
high-spin state.  It is a single Slater determinant, and every open
shell carries exactly one electron of the majority spin.  For such an
orbital the triple is
\begin{equation}
  n_{\alpha}=1,\qquad n_{\beta}=0,\qquad p_{2}=0,
  \label{eq:single-det-triple}
\end{equation}
so the weights are $w=(0,1,0,0)$ and
\begin{equation}
  s_{1}=0
  \label{eq:single-det-zero}
\end{equation}
for every open shell.  The nitrogen atom in its $^{4}S$ ground state
is the smallest example: three singly occupied $2p$ orbitals, and the
single-state criterion reports zero entropy for all three, although a
correct description of the atom requires all three in the active
space.

The zero is not a numerical accident.  In that determinant the
orbital is in a definite Fock state, so there is no entanglement
between it and the rest with respect to the spin degree of freedom.
What the determinant does not show is the degeneracy: the same
electronic structure appears in the $M_S=-S,\dots,+S$ components, and
the physical state is the manifold, not one projection of it.

\subsection{Merging the manifold}

Let $\{c^{(k)}\}_{k=1}^{R}$ be the CI vectors of the $R$ components,
with normalized columns $\sum_{I}|c^{(k)}_{I}|^{2}=1$, and let
$\omega_{k}$ be their ensemble weights.  The mixed one-orbital
reduced density matrix is the weighted mean of the per-state
matrices.  The map from the triple $(n_{\alpha},n_{\beta},p_{2})$ to
the weight vector in Eq.~\eqref{eq:four-state-weights} is affine, so
the weights of the mixture are the weighted mean of the per-state
weights:
\begin{equation}
  \bar{n}_{\alpha}=\sum_{k}\omega_{k}n^{(k)}_{\alpha},\qquad
  \bar{n}_{\beta}=\sum_{k}\omega_{k}n^{(k)}_{\beta},\qquad
  \bar{p}_{2}=\sum_{k}\omega_{k}p^{(k)}_{2}.
  \label{eq:manifold-mean}
\end{equation}
No density matrix is assembled, and the cost is one matrix product per
root: the per-state occupations follow from the determinant bitmaps
already used by the single-state route.

For a complete spin manifold the mean is spin-isotropic,
$\bar{n}_{\alpha}=\bar{n}_{\beta}$, and a singly occupied shell
obtains
\begin{equation}
  \bar{n}_{\alpha}=\bar{n}_{\beta}=\tfrac12,\qquad \bar{p}_{2}=0,
  \qquad w=\bigl(0,\tfrac12,\tfrac12,0\bigr),
  \qquad s_{1}=\ln 2 .
  \label{eq:manifold-ln2}
\end{equation}
The value $\ln 2$ is the entropy of one electron of unknown spin in
one orbital, which is exactly what the criterion should read for a
degenerate open shell.

\section{Which roots form the manifold}

For a pure spin state of spin $S$ the manifold has $2S+1$ components,
and in a calculation that returns the lowest eigenvalues in ascending
order they are the lowest $2S+1$ roots, because the components are
degenerate.  The gate therefore takes
\begin{equation}
  R = 2S+1,
  \label{eq:nroots}
\end{equation}
which the driver evaluates from the spin multiplicity of the request.
The assumption fails when a state of different spin and lower energy
lies inside the same window; the merged occupations then contain a
foreign state.  The worked example in
Sec.~\ref{sec:manifold-example} is the case where the assumption
holds and can be checked.

\section{API}

The analysis layer exposes the occupation triple and the entropy of a
manifold:
\begin{lstlisting}
from diffmrc.analysis.ci import (manifold_occupations,
                                 manifold_entropy_from_ci)

# vectors: shape (ndet,) or (ndet, nroots), in the determinant order
# produced by core.nevpt2.impl.cas_dets_for2
s1 = manifold_entropy_from_ci(vectors, dets, ncore=ncore, nact=nact)
n_a, n_b, p2 = manifold_occupations(vectors, dets, ncore=ncore,
                                    nact=nact)
\end{lstlisting}
Each column is normalized before use; the mixture weights carry the
ensemble, not the arbitrary scale of a vector.  The default is equal
weights.  A weight vector shorter than the root count is padded with
zeros, and negative weights, a zero weight sum, more weights than
states, and an all-zero column are refused with a message.

The differentiable sibling, used when the entropy is itself a function
of the orbital parameters, is
\texttt{driver.entropy.manifold\_entropy\_from\_ci\_masks} (masks
precomputed once per determinant list) and
\texttt{driver.entropy.manifold\_single\_orbital\_entropy\_from\_ci}
(determinant list given directly).  Both are bound to the NumPy route
by a test at $10^{-12}$.

The CI vectors of a converged calculation are obtained without
re-solving in a different basis:
\begin{lstlisting}
from diffmrc.driver.kernel import casci_vectors_at

vectors = casci_vectors_at(mol, C, space, theta, nroots=2 * S + 1)
\end{lstlisting}
\texttt{casci\_vectors\_at} shares the integral block and the orbital
rotation with the energy route, so the vectors belong to the reported
energy.

\section{The robust-level gate}
\label{sec:manifold-gate}

At the robust level the driver evaluates the merged spectrum after the
CASSCF has converged and records two entries of
\texttt{algorithm\_data}:
\begin{lstlisting}
{'active_space_verified': True | False,
 'active_space_entropy': {'s1': [...], 'nroots': R,
                          'verdict': '...', 'plateau_heights': [...],
                          'evidence': [...]}}
\end{lstlisting}
The check passes when the merged spectrum is not single-reference,
that is, when its largest single-orbital entropy reaches the
multireference criterion of $0.14$.  A verdict of
\texttt{single\_reference} fails the check, and the robust level
refuses the run: there is no active space to verify.  A verdict that
the spectrum is larger than the selected space is recorded but does
not veto the run, because the element rules add ligand orbitals whose
entropy lies below the line by construction; vetoing them would
contradict the selection rules themselves.

Two properties of the gate should be stated plainly.

\textbf{Scope.}  The spectrum is the one the selected space can show.
An orbital outside the space has zero entropy by construction, so the
gate can certify that the space is multireference and that its
orbitals are not noise; it cannot certify that the space is not
truncated at the frontier.  Completeness at the frontier remains the
responsibility of the selection rules.

\textbf{Failure to evaluate.}  If the spectrum cannot be evaluated the
gate reports unverified with the reason, rather than raising.  At the
robust level an unverified space is a refusal, which is the same
outcome the exception would have produced and a more useful record.

\section{Worked example: the nitrogen quartet}
\label{sec:manifold-example}

The $^{4}S$ ground state of the nitrogen atom, CAS(3,3) over the three
$2p$ orbitals with the 6-31G basis, at the ROHF reference orbitals.
The calculation returns four roots; Table~\ref{tab:manifold-n} lists
the single-orbital entropy of each root and of the merged manifold.

\begin{table}[h]
\centering
\caption{Single-orbital entropy of the N $^{4}S$ manifold,
CAS(3,3)/6-31G.  Each root is a component of the quartet; read
individually the roots disagree by three orders of magnitude, and the
merged value is $\ln 2$.}
\label{tab:manifold-n}
\begin{tabular}{lccc}
\toprule
root & $s_{1}(2p_x)$ & $s_{1}(2p_y)$ & $s_{1}(2p_z)$ \\
\midrule
0 & $6.60849\times10^{-1}$ & $6.60849\times10^{-1}$ & $6.60849\times10^{-1}$ \\
1 & $6.36514\times10^{-1}$ & $6.36514\times10^{-1}$ & $6.36514\times10^{-1}$ \\
2 & $2.2\times10^{-25}$ & $1.5\times10^{-24}$ & $5.1\times10^{-25}$ \\
3 & $1.68823\times10^{-1}$ & $1.68823\times10^{-1}$ & $1.68823\times10^{-1}$ \\
\midrule
merged & $6.93147\times10^{-1}$ & $6.93147\times10^{-1}$ & $6.93147\times10^{-1}$ \\
\bottomrule
\end{tabular}
\end{table}

Root 2 is the determinant of Eq.~\eqref{eq:single-det-zero} and
reports zero; roots 0, 1 and 3 are the other projections and report
values that depend on the projection.  The merged value equals
$\ln 2=0.693147$ for all three orbitals to the printed precision,
which is the spin-isotropic result of
Eq.~\eqref{eq:manifold-ln2} and confirms that the four roots form the
complete manifold.

For comparison, two closed-shell cases run through the same gate
report $4.787\times10^{-3}$ (H$_2$O/6-31G, CAS(2,2) over the two
frontier orbitals) and $6.792\times10^{-2}$ (H$_2$/STO-3G, CAS(2,2)).
Both lie below the $0.14$ criterion and are reported as
single-reference, which is the intended outcome: an active space is
not justified for them.

\section{Limitations}
\label{sec:manifold-limits}

\textbf{Root selection.}  Equation~\eqref{eq:nroots} assumes that the
lowest $2S+1$ roots are the components of the target manifold.  A
lower state of another spin inside the same window contaminates the
mixture.  The assumption can be checked by printing the root energies
and confirming degeneracy, which the driver records in the report but
does not enforce.

\textbf{Equal weights.}  The default weights are equal, which is the
state-averaged convention for a degenerate manifold.  A manifold split
by spin-orbit coupling is not degenerate; the API accepts explicit
weights, and the caller is responsible for choosing them.

\textbf{Within-space scope.}  As stated in
Sec.~\ref{sec:manifold-gate}, the spectrum cannot see orbitals that
the selection left out.

\textbf{A pre-existing failure on the quartet.}  The end-to-end
driver evaluates the orbital Hessian for its curvature check, and on
the N $^{4}S$ case that Hessian is not finite: at the converged point,
which sits within $2\times10^{-7}$ of the symmetric $\theta=0$, all
676 entries are NaN, while a random displacement of $0.01$ gives a
finite Hessian.  The exact degeneracy of the high-spin multiplet makes
the second derivative of the blocked eigensolver undefined.  The
driver now recovers the curvature from a symmetric difference of the
analytic gradient, which is finite at the same point, using a step of
$10^{-5}$ in the rotation parameters.  The route costs two gradient
evaluations per parameter and is capped at 300 parameters; a rotation
space with more than 300 parameters is refused with the diagnosis
rather than approximated by the difference route.  On N$_2$ the
spectra at steps $10^{-5}$, $10^{-4}$ and $10^{-3}$ agree (minimum
$0.0648$, maximum $66.96$, no negative eigenvalues), and
\texttt{algorithm\_data['curvature\_method']} records which route ran.
The numbers in Table~\ref{tab:manifold-n} were produced by calling
\texttt{casci\_vectors\_at} directly at the SCF orbitals.  The defect
is recorded here so that it is not mistaken for a property of the
merge.
